\section{Conclusion}

This study tests a narrow maintenance-warning question on simulated turbofan
engines: do measured channels add information beyond elapsed cycle and operating
conditions when the warning threshold is chosen without seeing the held-out
engines? The answer is yes for the tested N-CMAPSS subset. At a roughly 5\%
pre-window false-alert rate, all observable inputs raise final-window recall
from 55.4\% to 79.3\% over the cycle-plus-operating baseline. The paired
engine-cluster interval for the gain is 18.3 to 29.6 percentage points.

The telemetry-only result is also useful. With no cycle number or
operating-condition field, it recalls 77.1\% of final-window cycles and warns
69 of 90 engines by the window start. These findings are stable across five
random-forest seeds and use only the 14 documented measured channels. A
secondary policy that selects its model family inside the training folds raises
the all-observable result to AUROC 0.973 and recall 82.1\%, compared with 0.966
and 79.3\% for the fixed random forest at the same pooled false-alert rate.
This is a controlled performance gain, not a claim of a universally best model.

The contribution is evidence that observed sensor telemetry carries
maintenance-warning value beyond a flight-cycle clock in this simulator. It is
not evidence for a service interval or a fielded-aircraft decision rule. Those
claims require maintenance-defined costs and lead times, plus independent
labeled engine data.
